% fixture: koma/koma-sections
% tier: T3
% engine: document
% covers: KOMA-Script's sectioning: \section … \subsubsection in \sffamily\bfseries \Large, \large, \normalsize, ragged right, the number followed by \enskip, two \nobreak before the heading; run-in \paragraph and \subparagraph (indent 1em of the heading's font) after two \nobreak, toc lines with \nonumberline; starred headings reset the counters within them (a theorem's number); \label's name is the title
% issue: #360
\documentclass{scrartcl}
\usepackage{amsthm}
\newtheorem{theorem}{Theorem}[section]
\begin{document}
\section{Introduction}
Typesetting is the composition of text by means of arranging physical type or its digital equivalents.\label{sec:intro}
\begin{theorem}
A theorem in the first section.
\end{theorem}

\subsection{Background}
The breaking of paragraphs into lines is a small part of this craft, see Section~\ref{sec:intro}.

\subsubsection{A detail}
More text follows here.

\paragraph{Run in} A paragraph heading.

\subparagraph{Sub} A subparagraph heading.

\section*{Unnumbered}
Text after.
\begin{theorem}
Its number starts again within the section.
\end{theorem}
\section{A much longer section title that will need to break over two lines in the heading of the page}
Final words.
\end{document}
